We applied the method to the Dry Bean data of \citet{koklu2020drybean},
which contain 12 measured dimensions and four shape descriptors extracted from
images of 13,611 individual grains belonging to seven registered varieties.
The class labels were not used in estimation.  They were retained only to give
an external assessment of the fitted partition.  We standardized all 16
features and used the first two principal-component scores, thereby avoiding a
choice of variables based on the reported class labels.

Before running the bootstrap, we screened six data configurations (breast
cancer, wine, handwritten digits, and three pairs of bean varieties), two
dimensions, and Normal--Student-$t$ and Normal--skew-normal specifications.
The screen used numerical admissibility, component weights, effective kernel
support near the estimated boundary, and whether the boundary correction was
material.  It was used to select one regular illustration and one deliberately
difficult diagnostic case, rather than to select the largest apparent
improvement in clustering accuracy.  The complete screening results and failed
fits are retained in the reproduction archive.

For each selected data set we fitted a two-component Normal--Student-$t$
classification mixture.  We compared the analytic standard errors with 500
nonparametric bootstrap samples optimized locally within a coordinate-wise
trust region around the original solution.  This is the bootstrap analogue of
the local target in our theory.  We also performed 200 unrestricted refits from
the full fitting procedure.  The latter diagnostic allows the selected basin
to change and hence measures instability that is outside the local Gaussian
approximation.  Table~\ref{tab:applications} summarizes the results.

Figure~\ref{fig:application-boundaries} displays the standardized scores,
coloured by external variety labels, together with the fitted classification
boundary. Labels played no role in preprocessing or fitting.

\begin{figure*}[t]
\centering
\includegraphics[width=\textwidth]{figures/application_boundaries.pdf}
\caption{Dry Bean principal-component scores and fitted Normal--Student-$t$
classification boundaries. Colours show external labels used only for
evaluation; black curves are equality sets of the fitted weighted log densities.}
\label{fig:application-boundaries}
\end{figure*}

\begin{table}[t]
\centering
\scriptsize
\setlength{\tabcolsep}{3pt}
\caption{Dry Bean applications. The two final columns give the coordinate-
median ratio of an analytic standard error to the standard deviation from
successful local bootstrap fits. Boundary support is the number of observations
inside the fitted kernel tube. Local success is the percentage of 500 locally
optimized bootstrap refits that passed all convergence diagnostics.}
\label{tab:applications}
\resizebox{\linewidth}{!}{%
\begin{tabular}{lrrrrrr}
\toprule
Varieties & Sample size & Adjusted Rand & Kernel count & Local success & Naive ratio & Corrected ratio \\
\midrule
Dermason--Seker & 5573 & 0.867 & 33 & 95.0\% & 0.798 & 0.991 \\
Cali--Barbunya  & 2952 & 0.282 & 44 & 52.4\% & 0.538 & 0.850 \\
\bottomrule
\end{tabular}
}

\smallskip
\begin{minipage}{0.94\linewidth}\small
The median stabilized-to-naive standard-error inflation was 1.284 and 1.650,
respectively.  The percentages of unrestricted refits remaining within a
maximum coordinate distance of 0.75 from the original solution were 59.0\%
and 46.5\%, respectively.
\end{minipage}
\end{table}

\paragraph{A regular validation case.}
For Dermason--Seker the fitted partition agrees closely with the independently
recorded varieties (adjusted Rand index 0.867), both estimated component
weights are substantial, and 475 of 500 local bootstrap fits succeeded.  The
naive standard errors are systematically too small: their median ratio to the
local bootstrap standard deviation is 0.798.  Boundary-aware standard errors
raise this ratio to 0.991.  At the coordinate level the stabilized ratios range
from 0.886 to 1.082, whereas the naive ratios range from 0.576 to 0.939.  Thus
the aggregate improvement is not driven by one parameter.  The estimated
boundary curvature required no finite-sample shrinkage in this fit.

\begin{table}[t]
\centering
\scriptsize
\setlength{\tabcolsep}{3pt}
\caption{Dermason--Seker inference for interpretable fitted quantities. The
location distance is measured in standardized principal-component space.
Here ``s.e.'' means standard error and ``bootstrap s.d.'' means the standard
deviation of successful local-bootstrap estimates.}
\label{tab:dermason-interpretable}
\resizebox{\linewidth}{!}{%
\begin{tabular}{lrrrrl}
\toprule
Quantity & Estimate & Naive s.e. & Corrected s.e. & Bootstrap s.d. & Corrected 95\% interval\\
\midrule
Mixing weight $\pi_1$ & 0.345 & 0.006 & 0.008 & 0.008 & (0.329, 0.361)\\
$\|\boldsymbol\mu_1-\boldsymbol\mu_2\|$ & 5.202 & 0.036 & 0.037 & 0.035 & (5.130, 5.273)\\
Student-$t$ degrees of freedom $\nu$ & 20.680 & 3.848 & 4.184 & 5.533 & (12.479, 28.881)\\
\bottomrule
\end{tabular}
}
\end{table}

Scientifically, the correction changes uncertainty most clearly for the fitted
prevalence of the first hard-classification group: its standard error increases
by about 31\%, bringing it close to the local-bootstrap scale.  The estimated
location separation is so pronounced that its uncertainty changes little, as
predicted by the low-boundary-density argument.  For the Student-$t$ tail
parameter, the corrected interval is wider than the naive interval, although
the remaining gap to the local-bootstrap standard deviation warns that tail
estimation is harder.  The correction therefore changes substantive precision
selectively; it does not inflate every uncertainty by a common factor.

Figure~\ref{fig:application-calibration} confirms that the aggregate result is
not an artefact of taking medians. Dermason--Seker corrected ratios cluster
around one across the coordinate vector; Cali--Barbunya shows a substantial
but incomplete improvement for most coordinates.

\begin{figure*}[t]
\centering
\includegraphics[width=\textwidth]{figures/application_se_calibration.pdf}
\caption{Coordinate-level analytic standard errors divided by local-bootstrap
standard deviations. The dotted reference at one represents exact local
calibration.}
\label{fig:application-calibration}
\end{figure*}

Only 59.0\% of unrestricted bootstrap refits remained in the original local
basin.  Location-coordinate standard deviations from unrestricted refitting
were consequently orders of magnitude larger for some coordinates.  This does
not contradict the local calibration result: it identifies a separate layer
of optimization and selection uncertainty.  In practice, the local standard
errors should therefore be reported together with the basin-stability
diagnostic.

The unrestricted distances in Figure~\ref{fig:application-basins} make this
distinction visible. Both examples have a local group of refits, but
Dermason--Seker also has a separate group near distance five and
Cali--Barbunya has a broad nonlocal tail. A global standard deviation would
conflate these discrete changes with regular local sampling variation.

\begin{figure*}[t]
\centering
\includegraphics[width=\textwidth]{figures/application_basin_diagnostics.pdf}
\caption{Maximum coordinate distance between each unrestricted bootstrap
refit and the original solution. The dashed line is the prespecified local
trust-region radius 0.75.}
\label{fig:application-basins}
\end{figure*}

\paragraph{A nonregular diagnostic case.}
The Cali--Barbunya pair has stronger overlap, as reflected in the lower
external adjusted Rand index (0.282) and the larger median analytic correction
(a factor of 1.650).  The stabilized standard errors again move substantially
toward the local-bootstrap scale, improving the median ratio from 0.538 to
0.850.  However, only 262 of 500 local optimizations succeeded, and only 46.5\%
of unrestricted refits stayed near the original solution.  These diagnostics
fail the operational stability requirement needed to treat the local normal
approximation as a complete account of uncertainty.  We therefore present
this example as a stress test, not as evidence for routine Wald inference.
Together, the two applications show both the practical gain from estimating
boundary curvature and the need to distinguish local sampling uncertainty
from basin-selection instability.
